\chapter*{Introduction générale}
\thispagestyle{fancy}
\addcontentsline{toc}{chapter}{Introduction générale}

De nos jours, la santé mentale s'impose comme l'un des enjeux majeurs qui influencent profondément la qualité de vie et le bien-être de chacun d'entre nous. Selon l'Organisation mondiale de la santé~\cite{oms_mental_health}, près d'une personne sur huit vit avec un trouble mental, ce qui représente un défi considérable pour les systèmes de santé et la société. Pourtant, il reste difficile pour beaucoup d'accéder à un accompagnement psychologique personnalisé, en raison du manque de ressources spécialisées, de la stigmatisation persistante et du coût souvent élevé des soins.

Le diagnostic et le suivi des troubles psychiques reposent généralement sur de nombreux entretiens, questionnaires et évaluations cliniques. Ce processus peut s'avérer long, coûteux, et parfois source d'erreurs ou de retards dans la prise en charge. Face à ces défis, la recherche s'est tournée vers les nouvelles technologies, notamment l'intelligence artificielle, pour proposer des solutions capables d'analyser de grandes quantités de données et de fournir des recommandations personnalisées.

C'est dans ce contexte que s'inscrit le projet MindCare-AI, dont l'ambition est de proposer une approche innovante basée sur l'intelligence artificielle pour accompagner, évaluer et soutenir la santé mentale des utilisateurs. L'application mobile développée permet d'effectuer des auto-évaluations régulières, de suivre l'évolution du bien-être, d'obtenir des conseils personnalisés, d'interagir avec un chatbot intelligent, et d'analyser les signes vitaux de manière non-invasive par analyse d'image faciale.

Ce mémoire présente le développement complet de la solution MindCare-AI, depuis la conception jusqu'à l'implémentation et l'intégration des modules d'intelligence artificielle. Il est structuré en cinq chapitres :

\textbf{Le premier chapitre} présente le cadre général du projet en exposant la société d'accueil Mobelite, le contexte de la santé mentale digitale et la problématique à résoudre. Il analyse ensuite les solutions existantes sur le marché  à travers une étude comparative détaillée, avant de présenter notre solution innovante MindCare-AI. Le chapitre se termine par la présentation des méthodologies adoptées : l'approche Agile avec Scrum pour le développement logiciel et CRISP-DM pour les aspects d'intelligence artificielle.

\textbf{Le deuxième chapitre} se divise en deux parties complémentaires. La première partie \textit{(Initialisation du projet)} définit le périmètre fonctionnel en identifiant les acteurs du système , en spécifiant les besoins fonctionnels et non fonctionnels, en présentant l'écosystème technologique retenu , et en établissant le Product Backlog priorisé avec la planification stratégique des releases. La seconde partie \textit{(Analyse et spécification des besoins)} approfondit l'architecture technique en présentant le pattern MVC adopté, l'architecture globale du système avec ses composants principaux, et la conception détaillée formalisée par des diagrammes UML .

\textbf{Le troisième chapitre} détaille la Release 1 consacrée au développement du socle applicatif mobile à travers quatre sprints successifs. Le Sprint 1 implémente les fondations et l'authentification . Le Sprint 2 développe l'interface HomeScreen avec tableau de bord personnalisé. Le Sprint 3 crée l'assistant conversationnel empathique avec historique persistant .Enfin, Le Sprint 4 construit l'interface d'analyse rPPG  et dashboard de visualisations interactives avec filtrage temporel.

\textbf{Le quatrième chapitre} porte sur la Release 2 dédiée exclusivement à l'intégration ciblée de l'intelligence artificielle. En s'appuyant rigoureusement sur la méthodologie CRISP-DM, il présente le développement de deux modules d'IA complémentaires. Le Sprint 1 développe le module d'analyse rPPG des signes vitaux : collection et analyse du dataset PPG+DaLiA , préparation et prétraitement des signaux physiologiques , modélisation avec comparaison de plusieurs algorithmes , et intégration technique dans l'architecture existante avec validation de robustesse. Le Sprint 2 développe le chatbot émotionnel adaptatif : collection et analyse du dataset Emotions, traitement du déséquilibre des classes identifié (focus sur 4 émotions principales : joie, tristesse, colère, peur), développement d'un modèle CNN de deep learning atteignant 94.24\% de précision pour la classification émotionnelle textuelle, enrichissement du chatbot avec capacités d'adaptation contextuelle empathique, intégration du contexte physiologique (données rPPG) pour un accompagnement holistique multimodal, et validation conversationnelle de la pertinence et de l'empathie des réponses adaptées.

Ce travail démontre la faisabilité et la pertinence d'une approche intégrée combinant développement mobile moderne et intelligence artificielle appliquée à la santé mentale. Il ouvre ainsi de nouvelles perspectives pour l'accessibilité et la personnalisation des soins psychologiques, tout en positionnant MindCare-AI comme une solution innovante, scientifiquement validée et technologiquement avancée dans l'écosystème de la santé mentale digitale.